\documentclass[11pt, a4paper]{article}
\usepackage[a4paper, margin=2.5cm]{geometry}
\usepackage[T2A]{fontenc}
\usepackage[utf8]{inputenc}
\usepackage[english, russian]{babel}

\usepackage{enumitem}
\setlist[itemize]{label=-}
\usepackage{amsmath}
\usepackage{amssymb}
\usepackage{hyperref}
\title{\textbf{Ответы на вопросы для подготовки к лабораторной работе №3}\\ \large Методы измерений основных параметров усилителей высокочастотного диапазона}
\author{}
\date{}

\begin{document}

\maketitle

\begin{enumerate}[label=\textbf{\arabic*.}]

    \item \textbf{Методика измерения коэффициентов передачи и развязки с учетом ограничений аппаратуры}\\
    Для предотвращения выхода анализатора спектра из строя (максимальная входная мощность 1 Вт или +30 дБм) необходимо применять аттенюаторы и использовать калибровочные измерения .
    \begin{itemize}
        \item \textbf{Калибровка:} Генератор соединяется с анализатором спектра через аттенюатор 10 дБ и два измерительных кабеля, соединенных адаптером . На выбранных частотах измеряется уровень сигнала, что позволяет определить суммарные потери измерительного тракта ($IL$).
        \item \textbf{Измерение усиления ($S_{21}$):} Вместо адаптера подключается исследуемый усилитель. Для получения фактического коэффициента передачи из измеренного уровня на анализаторе вычитаются результаты калибровки.
        \item \textbf{Измерение развязки ($S_{12}$):} Усилитель подключается в обратном направлении (вход усилителя к анализатору, выход к генератору). Ослабление аттенюатора гарантирует, что вход анализатора спектра не перегрузится даже при случайном скачке сигнала.
    \end{itemize}

    \item \textbf{Расчет мощности шума на выходе широкополосного усилителя}\\
    Поскольку на вход подается шум с относительно высоким уровнем 0 дБм (что многократно превышает тепловой шум), собственным тепловым шумом усилителя можно пренебречь.
    \begin{itemize}
        \item $P_{\text{вых}} = P_{\text{вх}} + K_{p}$
        \item $P_{\text{вых}} = 0 \text{ дБм} + 10 \text{ дБ} = 10 \text{ дБм}$.
    \end{itemize}
    \textit{(Примечание: Если бы расчет проводился строго для теплового шума по формуле Найквиста, формула имела бы вид $P_{\text{вх}} + K_p + K_{\text{ш}}$, но в данном случае внешний шум доминирует).}

    \item \textbf{Отношение сигнал/шум на выходе усилителя}
    \begin{itemize}
        \item Спектральная плотность теплового шума при стандартной температуре: $G_0 = -174 \text{ дБм/Гц}$.
        \item Мощность теплового шума в полосе $\Pi = 10 \text{ МГц}$ ($10^7 \text{ Гц}$) на входе: \\ $P_{\text{ш.вх}} = -174 + 10\log_{10}(10^7) = -174 + 70 = -104 \text{ дБм}$.
        \item Отношение сигнал/шум на входе: $\text{SNR}_{\text{вх}} = P_{\text{сигн.вх}} - P_{\text{ш.вх}} = -80 \text{ дБм} - (-104 \text{ дБм}) = 24 \text{ дБ}$.
        \item Коэффициент шума усилителя показывает, на сколько децибел ухудшается отношение сигнал/шум при прохождении через устройство: $\text{SNR}_{\text{вых}} = \text{SNR}_{\text{вх}} - K_{\text{ш}}$.
        \item $\text{SNR}_{\text{вых}} = 24 \text{ дБ} - 6 \text{ дБ} = 18 \text{ дБ}$.
    \end{itemize}

    \item \textbf{Методика измерения коэффициента шума}\\
    Коэффициент шума измеряется косвенно путем сравнения уровня собственных шумов усилителя с тепловым шумом. При подключенном ко входу усилителя согласованном сопротивлении 50 Ом, измеряется уровень шума $P$ на анализаторе спектра. Коэффициент шума рассчитывается по формуле:
    \begin{equation*}
        K_{\text{ш}}[\text{дБ}] = P + IL_2 - (G_0 + \text{RBW}_{\text{в дБ}} + K_p) \quad
    \end{equation*}
    где $IL_2$ --- потери во втором кабеле, $G_0 = -174 \text{ дБм/Гц}$, а $\text{RBW}$ --- полоса пропускания фильтра анализатора в логарифмическом масштабе.

    \item \textbf{Настройки анализатора спектра для снижения уровня собственных шумов}\\
    Для понижения среднего отображаемого уровня шума (DANL) необходимо:
    \begin{itemize}
        \item Уменьшить полосу пропускания фильтра промежуточной частоты (уменьшить значение RBW).
        \item Включить встроенный предварительный усилитель (RF Preamp).
    \end{itemize}

    \item \textbf{Методика измерения точки децибельной компрессии ($P_{1\mathrm{dB}}$)}
    \begin{itemize}
        \item Точка компрессии характеризует границу линейного динамического диапазона. Изменяя уровень сигнала с генератора (например, от -10 до +20 дБм), строится зависимость выходной мощности от входной (амплитудная характеристика).
        \item По графику находится точка входной мощности, при которой коэффициент усиления становится на 1 дБ меньше, чем на идеальном линейном участке. Для точности предварительно вычитаются суммарные потери измерительного тракта ($IL_2$).
    \end{itemize}

    \item \textbf{Изменение точки децибельной компрессии с ростом частоты}\\
    С ростом частоты коэффициент усиления большинства широкополосных усилителей начинает падать из-за влияния паразитных емкостей транзисторов. Как следствие, максимальная выходная мощность и точка децибельной компрессии $P_{1\mathrm{dB}}$ снижаются.

    \item \textbf{Методика измерения $IP_3$ и выбор между смесителем и сумматором}\\
    Точка пересечения 3-го порядка ($IP_3$) измеряется путем подачи на вход усилителя двухтонового сигнала и оценки отношения мощностей основных тонов к продуктам интермодуляции на выходе.
    \begin{itemize}
        \item Использовать делитель/сумматор мощности предпочтительнее, так как он является линейным пассивным элементом и обеспечивает высокую изоляцию генераторов друг от друга. Смеситель является принципиально нелинейным устройством и сам генерирует собственные комбинационные составляющие, которые наложатся на интермодуляционные искажения усилителя и приведут к ошибкам измерения.
    \end{itemize}

    \item \textbf{Частоты двухтонового сигнала и продуктов интермодуляции}
    \begin{itemize}
        \item Для формирования разницы в 1 МГц выберем тоны: $f_1 = 1000 \text{ МГц}$ и $f_2 = 1001 \text{ МГц}$.
        \item Продукты интермодуляции 3-го порядка возникают на частотах $2f_1 - f_2$ и $2f_2 - f_1$.
        \item Частоты продуктов: $2(1000) - 1001 = 999 \text{ МГц}$ и $2(1001) - 1000 = 1002 \text{ МГц}$.
    \end{itemize}

    \item \textbf{Изменение $IP_3$ с ростом частоты}\\
    Поскольку с ростом частоты падают усилительные свойства активных компонентов и увеличивается влияние нелинейных паразитных реактивностей, линейность усилителя ухудшается. Следовательно, значения $IIP_3$ и $OIP_3$ с ростом частоты уменьшаются.

    \item \textbf{Влияние аттенюатора 6 дБ на входе на параметры $IIP_3$ и $OIP_3$}
    \begin{itemize}
        \item Выходная точка пересечения $OIP_3$ является характеристикой самого усилителя и \textbf{не изменится}, так как нелинейный элемент генерирует те же искажения при достижении того же уровня на выходе.
        \item Входная точка пересечения $IIP_3$ системы \textbf{увеличится на 6 дБ} ($IIP_{3\text{.sys}} = IIP_3 + 6 \text{ дБ}$). Чтобы получить те же искажения на выходе усилителя, на вход схемы (перед аттенюатором) теперь нужно подать сигнал на 6 дБ мощнее.
    \end{itemize}

    \item \textbf{Ширина полосы продуктов интермодуляции 3-го порядка}\\
    Кубическая нелинейность в спектральной области приводит к трехкратной свертке спектра сигнала самого с собой. Поэтому, если исходный сигнал занимает полосу 10 МГц, спектр продуктов интермодуляции 3-го порядка расширится в 3 раза и займет полосу \textbf{30 МГц} на той же центральной частоте 1 ГГц.

    \item \textbf{Оценка изменения коэффициента шума при перестановке аттенюатора}\\
    Коэффициент шума пассивного аттенюатора численно равен его подавлению ($K_{\text{ш.атт}} = 6 \text{ дБ}$, в разах $F_{\text{атт}} \approx 3.98$). Согласно формуле Фрииса для многокаскадного тракта: $K_{\text{ш}\Sigma} = K_{\text{ш}1} + \frac{K_{\text{ш}2} - 1}{K_{p1}}$.
    \begin{itemize}
        \item \textbf{Случай 1 (аттенюатор на входе):} Весь шум аттенюатора складывается с шумом усилителя. \\ $K_{\text{ш}\Sigma}[\text{дБ}] = L_{\text{атт}} + K_{\text{ш.ус}} = 6 \text{ дБ} + 3 \text{ дБ} = \textbf{9 дБ}$.
        \item \textbf{Случай 2 (аттенюатор на выходе):} Собственный шум второго каскада (аттенюатора) делится на коэффициент усиления первого каскада ($K_{p} = 20 \text{ дБ}$, в разах $G = 100$). 
        \begin{equation*}
            F_{\Sigma} = F_{\text{ус}} + \frac{F_{\text{атт}} - 1}{G_{\text{ус}}} = 10^{0.3} + \frac{3.98 - 1}{100} = 1.995 + 0.0298 = 2.0248
        \end{equation*}
        В децибелах: $K_{\text{ш}\Sigma}[\text{дБ}] = 10\log_{10}(2.0248) \approx \textbf{3.06 дБ}$.
    \end{itemize}
    \textbf{Вывод:} Коэффициент шума системы \textbf{улучшится (уменьшится) почти на 6 дБ} при переносе аттенюатора со входа на выход.

\end{enumerate}

\end{document}